% Part: lambda-calculus
% Chapter: introduction
% Section: overview

\documentclass[../../../include/open-logic-section]{subfiles}

\begin{document}

\olfileid{lam}{int}{ovr}
\olsection{आढावा}

अॅलॉन्झो चर्च यांनी १९३०च्या दशकाच्या सुरुवातीला लॅम्डा कलनाची रचना
रचनाशील तर्कशास्त्राचा आधार म्हणून केली होती; संगणनक्षम फलनांचे
प्रतिमान म्हणून \emph{नव्हे}. पण ते ट्यूरिंग-संगणनक्षम फलने आणि आंशिक
पुनरावर्ती फलने यांसारख्या संगणनक्षमतेच्या इतर व्याख्यांशी समतुल्य
असल्याचे लवकरच दाखवले गेले. सुरुवातीला हे किंचित अनपेक्षित वाटले,
ही गोष्ट हे वैशिष्ट्यीकरण अधिकच रंजक बनवते.

फलनाला स्वतंत्र नाव देण्याऐवजी, त्याची व्याख्या करणाऱ्या सांकेतिक
व्यक्तीकरणाने थेट त्याचा निर्देश करण्याची लॅम्डा संकेतन ही सोयीची
पद्धत आहे. “$f$ हे $f(x) = x + 3$ ने परिभाषित फलन असू द्या,” असे
म्हणण्याऐवजी “$f$ हे $\lambd[x][(x + 3)]$ फलन असू द्या,” असे म्हणता
येते. दुसऱ्या शब्दांत, $\lambd[x][(x+3)]$ हे फलसाधकाच्या मूल्यात तीन मिळवणाऱ्या
फलनाचे केवळ \emph{नाव} आहे. या व्यक्तीकरणातील $x$ हे स्थानधारक आहे: त्याच फलनाचा निर्देश
$\lambd[y][(y + 3)]$ नेही करता येतो. इतर प्राचले असतानाही हे संकेतन
चालते. उदाहरणार्थ, $g(x, y)$ हे दोन चरांचे फलन आणि $k$ ही नैसर्गिक
संख्या असेल, तर $\lambd[x][g(x,k)]$ हे कोणत्याही~$x$ ला
~$g(x, k)$ शी जोडणारे फलन आहे.

सांकेतिक व्यक्तीकरणावरून फलन परिभाषित करण्याच्या या पद्धतीला
\emph{लॅम्डा अमूर्तीकरण} म्हणतात. लॅम्डा अमूर्तीकरणाची दुसरी बाजू
म्हणजे \emph{उपयोजन}: आपल्याकडे एखादे फलन $f$ असेल (समजा, नैसर्गिक
संख्यांवर परिभाषित), तर त्याचे 2 सारख्या कोणत्याही मूल्यावर
\emph{उपयोजन} करता येते. अर्थात, प्रचलित संकेतनात आपण मिळणाऱ्या
मूल्यासाठी~$f(2)$ लिहितो.

लॅम्डा अमूर्तीकरण आणि उपयोजन एकत्र केल्यास काय होते? मग अमूर्तीकृत
चराच्या जागी ज्यावर उपयोजन केले ते मूल्य “बसवून” परिणामी
व्यक्तीकरण सोपे करता येते. उदाहरणार्थ,
\[
(\lambd[x][(x + 3)])(2)
\]
हे~$2 + 3$ असे सोपे करता येते.

आतापर्यंत आपण प्रचलित संकल्पनांसाठी नवे संकेतन मांडण्यापलीकडे काही
केलेले नाही. मात्र लॅम्डा कलन संचसैद्धान्तिक दृष्टिकोनापासून अधिक
मूलगामी फारकत घेते. या चौकटीत:
\begin{enumerate}
\item सर्वकाही फलन दर्शवते.
\item लॅम्डा अमूर्तीकरणाने फलने परिभाषित करता येतात.
\item कोणत्याही गोष्टीचे इतर कोणत्याही गोष्टीवर उपयोजन करता येते.
\end{enumerate}
उदाहरणार्थ, $F$ हे लॅम्डा कलनातील पद असेल, तर $F(F)$ अर्थपूर्ण
असल्याचे नेहमी गृहीत धरले जाते. या उदार चौकटीला
\emph{प्रकारविरहित} लॅम्डा कलन म्हणतात; येथे “प्रकारविरहित” याचा अर्थ
“कशाचे कशावर उपयोजन करता येईल यावर बंधन नाही” असा आहे.

\begin{digress}
\emph{प्रकारयुक्त} लॅम्डा कलनही आहे; तो प्रकारविरहित आवृत्तीचा
एक महत्त्वाचा पर्याय आहे. प्रकारयुक्त लॅम्डा कलन अनेक बाबतींत
प्रकारविरहित कलनासारखे असले, तरी त्याची अभिजात संचसैद्धान्तिक
चौकटीशी सांगड घालणे खूप सोपे आहे आणि त्याचे काही गुणधर्मही
खूप वेगळे आहेत.

लॅम्डा कलनावरील संशोधन सैद्धान्तिक संगणकशास्त्रात आणि
प्रोग्रामिंग भाषांच्या रचनेत मध्यवर्ती ठरले आहे. जॉन मॅकार्थी
यांनी १९५०च्या दशकात तयार केलेली LISP ही या कल्पनांचा
प्रभाव असलेल्या भाषेचे सुरुवातीचे उदाहरण आहे.
\end{digress}

\end{document}
